\section*{Chapitre \ref{ch:g12:climate-change} --- Le changement climatique~: observations, causes et projections}

\begin{solution}{exo:g12:climate-change:1}
La variation qu'il provoque du flux net d'énergie qui entre dans le système
climatique au sommet de l'atmosphère, avant que la température de surface
réagisse~; on le mesure en \unit{W/m^2}. Positif~: la planète gagne de
l'énergie et se réchauffe~; négatif~: elle perd de l'énergie et se
refroidit.
\end{solution}

\begin{solution}{exo:g12:climate-change:2}
La plus haute en mai, la plus basse en septembre. Au printemps et en été,
la végétation de l'hémisphère Nord, où pousse la plupart des plantes
terrestres, absorbe par photosynthèse plus de dioxyde de carbone que les
plantes et les sols n'en libèrent~; en automne et en hiver, la
décomposition et la respiration dominent, et la concentration remonte.
\end{solution}

\begin{solution}{exo:g12:climate-change:3}
\emph{Très probable}~: une probabilité de 90 à 100\,\%~; \emph{probable}~:
66 à 100\,\%. $[0.95 \text{ à } 1.20]$ est une fourchette très probable~: il
y a 90\,\% de chances que la vraie valeur se trouve entre 0.95 et 1.20.
\end{solution}

\begin{solution}{exo:g12:climate-change:4}
91\,\%. L'eau se dilate en se réchauffant~: un océan plus chaud occupe un
plus grand volume, et sa surface monte même si aucune eau n'est ajoutée (la
fonte des glaces continentales ajoute de l'eau en plus).
\end{solution}

\begin{solution}{exo:g12:climate-change:5}
La vapeur d'eau~: un air plus chaud contient davantage de vapeur d'eau, un
gaz à effet de serre, ce qui renforce le réchauffement. La neige et la
glace~: leur fonte découvre des surfaces plus sombres, qui absorbent
davantage de lumière solaire. La réponse négative~: une planète plus chaude
rayonne davantage d'infrarouge vers l'espace, ce qui la ramène à
l'équilibre.
\end{solution}

\begin{solution}{exo:g12:climate-change:6}
$427.35 - 315.98 = \qty{111.37}{ppm}$~; $111.37/315.98 \approx 0.352$, soit
une hausse d'environ 35\,\%. Sur 66 ans~: $111.37/66 \approx
\qty{1.7}{ppm}$ par an en moyenne.
\end{solution}

\begin{solution}{exo:g12:climate-change:7}
$F = 3.93 \times \ln(427.35/278)/\ln 2 = 3.93 \times 0.4300/0.6931 \approx
\qty{2.44}{W/m^2}$.
\end{solution}

\begin{solution}{exo:g12:climate-change:8}
$3.7 \times 30 = \qty{111}{mm}$, environ \qty{11}{cm}. $3.7/1.3 \approx 2.8$~:
presque trois fois plus vite.
\end{solution}

\begin{solution}{exo:g12:climate-change:9}
$\Delta T = 2.72/1.16 \approx \qty{2.3}{\degreeCelsius}$. Le climat n'est
pas à l'équilibre~: l'océan, qui absorbe environ 90\,\% de l'énergie
supplémentaire, se réchauffe lentement, si bien que la surface n'a pas
encore rattrapé le \omterm{def:g12:climate-change:forcing}{forçage} (le système gagne encore environ
\qty{0.8}{W/m^2}).
\end{solution}

\begin{solution}{exo:g12:climate-change:10}
Le Soleil et les volcans expliquent au plus \qty{0.1}{\degreeCelsius} d'un
réchauffement de \qty{1.06}{\degreeCelsius}~: ils ne peuvent pas en être la
cause. L'autre cause naturelle est la variabilité interne du climat (ses
fluctuations naturelles, comme les échanges de chaleur entre l'océan et
l'atmosphère)~; l'évaluation fixe sa contribution à $-0.2$ à $+\qty{0.2}{\degreeCelsius}$,
elle aussi bien trop faible.
\end{solution}

\begin{solution}{exo:g12:climate-change:11}
Très probablement sous \qty{2}{\degreeCelsius}~: seulement SSP1-1.9
(fourchette très probable de 1.0 à 1.8). Très probablement au-dessus~:
SSP2-4.5 (2.1 à 3.5), SSP3-7.0 (2.8 à 4.6) et SSP5-8.5 (3.3 à 5.7). Pour
SSP1-2.6 (1.3 à 2.4), aucun des deux n'est très probable.
\end{solution}

\begin{solution}{exo:g12:climate-change:12}
Le \omterm{def:g12:climate-change:forcing}{forçage} des aérosols vaut $-0.22 - 0.84 = \qty{-1.06}{W/m^2}$~; sans lui,
le total devient $2.72 + 1.06 = \qty{3.78}{W/m^2}$. Le réchauffement à
l'équilibre augmente de $1.06/1.16 \approx \qty{0.9}{\degreeCelsius}$.
\end{solution}

\begin{solution}{exo:g12:climate-change:13}
Comme $\ln(2C) - \ln C = \ln 2$ quel que soit $C$,
\[
  F(556) - F(278) = \frac{3.93\,\ln 2}{\ln 2} = \qty{3.93}{W/m^2}, \qquad
  F(1112) - F(556) = \frac{3.93\,\ln 2}{\ln 2} = \qty{3.93}{W/m^2}:
\]
le même \omterm{def:g12:climate-change:forcing}{forçage} pour chaque doublement. Total à \qty{1112}{ppm}~: $2 \times 3.93 =
\qty{7.86}{W/m^2}$.
\end{solution}

\begin{solution}{exo:g12:climate-change:14}
$\dfrac{\dd F}{\dd C} = \dfrac{F_{2\times}}{C \ln 2}$. À \qty{280}{ppm}~:
$3.93/(280 \times 0.6931) \approx \qty{0.0202}{W/m^2}$ par ppm~; à
\qty{420}{ppm}~: $\approx \qty{0.0135}{W/m^2}$ par ppm. Chaque ppm
supplémentaire ajoute moins de \omterm{def:g12:climate-change:forcing}{forçage} à mesure que la concentration
augmente (1.5 fois moins à 420 qu'à 280), mais la concentration continue
de monter, si bien que le \omterm{def:g12:climate-change:forcing}{forçage} augmente encore.
\end{solution}

\begin{solution}{exo:g12:climate-change:15}
$3.93/3.22 \approx \qty{1.2}{\degreeCelsius}$ sans \omterm{def:g12:climate-change:feedback}{rétroactions}~;
$3.93/1.16 \approx \qty{3.4}{\degreeCelsius}$ avec elles. Facteur
d'amplification~: $3.22/1.16 \approx 2.8$.
\end{solution}

\begin{solution}{pb:g12:climate-change:1}
\textbf{1.} $3.37 + 3.51 = \qty{6.88}{ppm}$, environ \qty{7}{ppm}.

\textbf{2.} Du printemps à la fin de l'été, les plantes de l'hémisphère
Nord retirent le dioxyde de carbone de l'air par photosynthèse plus vite
que la respiration et la décomposition ne le restituent~: la concentration
baisse jusqu'en septembre. En automne et en hiver, le bilan s'inverse, et
la concentration remonte jusqu'à son maximum de mai.

\textbf{3.} L'hémisphère Sud a beaucoup moins de terres, et donc beaucoup
moins de végétation, que l'hémisphère Nord~; l'essentiel de sa surface est
océanique.

\textbf{4.} En altitude, il se trouve au-dessus de la couche d'air pollué
proche du sol~; loin des continents et des villes, l'air qu'il
échantillonne est bien mélangé et représente l'atmosphère entière, non une
source locale.

\textbf{5.} $(325.68 - 316.91)/10 \approx \qty{0.88}{ppm}$ par an.

\textbf{6.} $(427.35 - 401.01)/10 \approx \qty{2.63}{ppm}$ par an.

\textbf{7.} $2.634/0.877 \approx 3.0$~: le rythme de hausse a triplé.

\textbf{8.} $(427.35 - 278)/278 \approx 0.537$~: environ 54\,\%.

\textbf{9.} $427.35 - 300 \approx \qty{127}{ppm}$.

\textbf{10.} $F = 3.93 \times \ln(427.35/278)/\ln 2 \approx
\qty{2.44}{W/m^2}$.

\textbf{11.} $F(316.91) = 3.93 \times \ln(316.91/278)/\ln 2 \approx
\qty{0.74}{W/m^2}$~; augmentation~: $2.44 - 0.74 \approx \qty{1.70}{W/m^2}$.

\textbf{12.} Un \unit{W/m^2} de plus multiplie la concentration par
$2^{1/3.93} \approx 1.193$~: $427.35 \times 1.193 \approx \qty{510}{ppm}$.

\textbf{13.} $(509.8 - 427.35)/2.634 \approx 31$ ans.

\textbf{14.} $2 \times 278 = \qty{556}{ppm}$~; \omterm{def:g12:climate-change:forcing}{forçage}~: $F_{2\times} =
\qty{3.93}{W/m^2}$.

\textbf{15.} Par définition de la sensibilité, \qty{3}{\degreeCelsius}
(probablement de 2.5 à \qty{4}{\degreeCelsius}).

\textbf{16.} $\lambda = 3.93/3 \approx \qty{1.31}{W.m^{-2}.\degreeCelsius^{-1}}$.
La somme des \omterm{def:g12:climate-change:feedback}{rétroactions}, \qty{1.16}{W.m^{-2}.\degreeCelsius^{-1}}, donne
$3.93/1.16 \approx \qty{3.4}{\degreeCelsius}$, un peu plus~; la meilleure
estimation évaluée combine cette physique avec le réchauffement observé et
les climats passés, et les deux valeurs se situent dans la fourchette
probable.

\textbf{17.} L'océan stocke l'essentiel de l'énergie supplémentaire et se
réchauffe lentement~: la surface n'atteint son réchauffement à l'équilibre
qu'au bout de décennies, voire de siècles.

\textbf{18.} Il suppose un récit des émissions futures (la hausse des dix
dernières années maintenue pour toujours)~: c'est un \omterm{def:g12:climate-change:scenario}{scénario}. Une
\omterm{def:g12:climate-change:scenario}{projection} calcule la réponse du climat et du cycle du carbone à un
\omterm{def:g12:climate-change:scenario}{scénario d'émissions} avec des modèles physiques, et donne des fourchettes
et des probabilités.

\textbf{19.} $(556 - 427.35)/2.634 \approx 48.8$ ans après 2025~: vers
\textbf{2074}.
\end{solution}
